% ============================================================================
% Reference-points subsection + table — DROP-IN for experiments.tex
% Status: DRAFT — placeholders {DL_XX}/{PT_XX} to be filled by
% plan/gates/gate6_baselines/integrate_gate6.py after gate6_audit.py passes.
% Framing: anchors, not peers (user decision 2026-09-01).
% ============================================================================

\subsection{External reference points}
\label{sec:refpoints}
The comparisons above are internal to the family: every row shares one
backbone and differs only in the attention module. To make the absolute
MSE levels interpretable, we additionally ran two widely used published
forecasting architectures under the \emph{identical frozen protocol} (task,
split, train-only normalization, optimizer, schedule, five paired seeds) as
\emph{reference points}, not as peer methods: \textbf{DLinear}
\citep{zeng2023dlinear} (channel-shared configuration) and \textbf{PatchTST}
\citep{nie2023patchtst} (official supervised ETTh1 configuration: patch 16,
stride 8, $d_{\mathrm{model}}=16$, 3 layers, dropout 0.3, RevIN). Their edge
over our minimal single-block backbone comes from their full architecture
(patching, channel independence, decomposition; in DLinear's case no
attention at all), not from their attention; the comparison is therefore
confounded by design, and \emph{no competitiveness claim is made or
implied}. The purpose is calibration: what an interested reader needs in
order to locate the family's absolute numbers. Because the protocol is held
fixed while the architectures are not peers, the anchors are
protocol-comparable but not method-comparable.

Table~\ref{tab:refpoints} reports the result. The anchors are decisively
ahead on the real-world sets---most starkly, DLinear (49K parameters, no
attention) outperforms every configuration of the family on ETTh1 validation
and test. This is the LTSF-linear result of \citet{zeng2023dlinear}
reproduced under our frozen protocol, and it re-anchors the paper's reading:
at this scale, the attention module is not the accuracy bottleneck, and the
design-space structure measured in RQ1--RQ3 is a property of the attention
family itself, orthogonal to backbone competitiveness. Consistent with this,
the internal ordering of the family (compressed $>$ full attention on most
datasets; decoupled-rate effects) is a statement within the shared backbone
and is unchanged by the anchors' presence. PatchTST is reported only where
its channel-independent training fits the single-machine CPU budget (the
ETT family, exchange, and the synthetic probes); on the wide datasets
(solar 137, electricity 321, traffic 862 channels) it is a \emph{declared
omission} (${\sim}3$--$13$\,h per seed on CPU), not a missing result.

\begin{table}[t]
\centering
\small
\caption{External reference points under the frozen protocol (5 paired
seeds; validation MSE in standardized units, lower is better). Reference
points, not peer methods: the anchors carry their full published
architectures (DLinear has no attention) while Normal/CSA/DRRA share one
minimal backbone (rows from Table~\ref{tab:coverage}). PatchTST is run
where it fits the CPU budget and is a declared omission on the wide
datasets (solar/electricity/traffic; ---). Bold: best reference point per
row.}
\label{tab:refpoints}
\setlength{\tabcolsep}{4pt}
\resizebox{\linewidth}{!}{%
\begin{tabular}{lccccc}
\toprule
Dataset & DLinear & PatchTST & Normal & CSA $(4,4)$ & DRRA $(8,4)$ \\
\midrule
% TABLE_BODY — generated by integrate_gate6.py (family columns verified
% bit-identical to tab:coverage; anchor columns from gate6 audit)
\bottomrule
\end{tabular}}
\end{table}
